\documentclass[12pt,twoside,letterpaper]{article}
\usepackage[T1]{fontenc}
\usepackage{tabularx}
\usepackage{array}

%NOTE: This report format is 

\newcommand{\reporttitle}{{\huge \textsc{Rapport} \\ \vspace{0.5cm} {\Large Travail pratique 1}}}

\newcommand{\reportauthorOne}{Pier-Luc Dupéré}
\newcommand{\reportProfessor}{Yacine Yaddaden, Ph. D.}
\newcommand{\reporttype}{Coursework}
\bibliographystyle{plain}

% include files that load packages and define macros
\input{includes} % various packages needed for maths etc.


%%%%%%%%%%%%%%%%%%%%%%%%%%%%

\begin{document}
% front page
\input{titlepage}


%%%%%%%%%%%%%%%%%%%%%%%%%%% table of content
%If a table of content is needed, simply uncomment the following lines
\tableofcontents
\newpage

%%%%%%%%%%%%%%%%%%%%%%%%%%%% Main document
\section{Introduction}
pld

\section{Choix, exploration et compréhension des données}
Le projet OASIS[] met gratuitement à la disposition des chercheurs des données d'imagerie du 
cerveau. On a utiliser la base de données OASIS 1 qui contient des images à résonance magnétique
ayant été diagnostiquer pas des médecins professionnels. Dans cette base de données, on retrouve 
le résultats des patients et qui sont classifier en 4 catégories. De plus, la base de données contient 
d'autres informations et évaluations, comme les résultats de tests cognitifs et l'âge des patients etc.

\newpage
\section{Prétraitement et représentation des données}
En ce qui concerne la préparation des données, on a extrait des coupes axial du cerveau de chaque 
patients. Pour faire la sélection des tranches, on a utiliser un logiciel appeler XMedCon[]. On a
pris 25 coupes par échantillon. Les coupes de 80 à 104 inclusivement de chaque fichier hdr.
\newline

Dans ce travail, seulement 15 échantillons pour chacune des trois catégories : 
\begin{itemize}
    \item Démence légère (Mild Demented)
    \item Sans démence (Non Demented)
    \item Démence très légère (Very Mild Demented)
\end{itemize}

Le code pré-traite une image IRM du cerveau. Il la convertit en niveaux de gris, applique un 
flou gaussien, puis utilise le seuillage d:Otsu et des opérations morphologiques pour nettoyage
de l'image. Il repère ensuite les contours situés dans la zone centrale afin de créer un masque 
du cerveau. Ce masque est appliqué à l'image, puis la zone obtenue est recadrée et redimensionnée 
en 128 × 128 pixels pour faciliter son analyse. On applique cela a tout les IRM et on sauvegarde le 
tout sous forme de png. 

\begin{figure}[htbp]
    \centering
    \includegraphics[width=0.95\textwidth]{img/pretraitement.JPG}
    \caption{Étapes du prétraitement des images IRM}
    \label{fig:pretraitement}
\end{figure}


\newpage
\section{Sélection / réduction de dimension}
Le Local Binary Pattern (LBP) de la librairie scikit-image[] est utilisé afin de réduire et 
sélectionner l'information des IRM. Cette fonction compare les pixels voisins au pixel central. 
Cela forment un motif binaire qui représente la texture autour du pixel. On peut ensuite utiliser 
les motifs obtenus pour générer un histogramme pour caractériser et classer les textures.
Dans notre étude, la classification des textures revient à classer les échantillons selon leur catégorie.
\newline

Les parameters de la fonction sont les suivants:
\begin{itemize}
    \item P=8 (nombre de point)
    \item R=3 (rayon)
    \item method="default" (changements d'intensité en niveaux de gris sans rotations.)
\end{itemize}

\begin{figure}[htbp]
    \centering
    \includegraphics[width=0.95\textwidth]{img/LBP.JPG}
    \caption{À gauche : représentation de l'image LBP ; à droite : histogramme correspondant}
    \label{fig:LBP}
\end{figure}


De plus, une réduction de dimension par analyse en composantes principales de type PCA est appliquée 
aux données dans le pipeline[] d'execution. 

Les parameters du PCA:
\begin{itemize}
    \item \texttt{n\_components=0.95} (\% de la variance des données conservée)
    \item \texttt{svd\_solver=auto} (méthode de calcul utilisé)
    \item \texttt{power\_iteration\_normalizer=auto}
\end{itemize}


\newpage
\section{Compréhension et recherche sur l'algorithme}


\begin{figure}[htbp]
    \centering
    \includegraphics[width=0.95\textwidth]{img/SVM.JPG}
    \caption{Pipeline réalise une classification avec un SVM}
    \label{fig:SVM}
\end{figure}

Dans l'image ci-dessus, on peut voir un pipeline qui standardise les caractéristiques et puis qui réduit 
leur nombre en conservant 95 \% de la variance. Par la suite, on utilise un SVM pour classer les échantillons.

Dans le tableau ci-dessous, il y a un résumé synthèse de tout les principale hyperparamètres. On peut voir 
la fonction que chacun de chacun ainsi que leur effets sur l'hyperplan qui segmente les échantillons. De plus, 
les valeurs typique qui vont être tester plus tard.

\begin{table}[htbp]
    \centering
    \caption{Principaux hyperparamètres d'un SVM []}
    \label{tab:hyperparametres-svm}
    \begin{tabularx}{\textwidth}{
        |>{\ttfamily}l
        |X
        |X
        |>{\ttfamily}l|
    }
        \hline
        \normalfont\textbf{Hyperparamètre} &
        \textbf{Fonction} &
        \textbf{Effets} &
        \normalfont\textbf{Valeurs} \\
        \hline
        kernel &
        Type de noyau utilisé &
        Détermine la forme de la frontière de décision &
        \normalfont\texttt{linear}, \texttt{poly}, \texttt{rbf}, \texttt{sigmoid} \\
        \hline

        C &
        Pénalisation des erreurs de classification &
        Une faible valeur favorise une marge plus large et peut entraîner du sous-apprentissage ; une valeur élevée pénalise davantage les erreurs et peut favoriser le surapprentissage &
        \normalfont{0.01 à 100} \\
        \hline

        gamma &
        Inverse du rayon d'influence, notamment pour le noyau RBF &
        Une faible valeur produit une frontière plus lisse ; une valeur élevée peut la rendre plus complexe &
        \normalfont{0.001 à 10.0} \\
        \hline

        degree &
        Degré du polynôme utilisé par le noyau polynomial &
        Un degré plus élevé permet une frontière de décision plus flexible &
        \normalfont{2 à 5} \\
        \hline

        coef0 &
        Terme constant des noyaux polynomial et sigmoid &
        Influence la forme du noyau, notamment dans le cas du noyau polynomial &
        \normalfont{0.0 à 1.0} \\
        \hline
    \end{tabularx}
\end{table}


\newpage
\section{Protocole de validation et métriques}
Afin de cherche la meilleure combinaison de paramètres pour un modèle. On a utiliser la librairie 
GridSearchCV[]. Elle teste les combinaisons indiquées et évalue chacune avec une validation croisée, 
puis sélectionne celle qui obtient le meilleur score.

\begin{figure}[htbp]
    \centering
    \includegraphics[width=0.95\textwidth]{img/gridsearchcv.JPG}
    \caption{Utilisation de GridSearchCV}
    \label{fig:GridSearchCV}
\end{figure}


Pour chaque combinaison, GridSearchCV évalue les performances en validation croisée (cv), avec l'accuracy 
comme mesure. Il teste chacune des combinaisons répétée pour les différentes validations croisées avec 
StratifiedKFold[]. Cela signifie que chaque combinaison de paramètres sera évaluée sur ces 10 découpages, 
puis comparée selon son accuracy moyenne.

Après l'entrainment du model, on prend le meilleur résultat avec les meilleurs paramètres. 


Pour finir, le meilleur modèle est évalué sur les données de test, qu'il n'a pas vues pendant 
son entraînement. Il prédit la classe de chaque échantillon. Ces prédictions sont comparées 
aux vraies classes pour calculer son exactitude et produire des documents afin de faciliter 
l'analyse des performances du modèle sur les données de test.
\newpage
\section{Présentation et interprétation des résultats}




\newpage
\section{Bilan critique de l'algorithme}
Les méthodes à noyaux offrent une grande flexibilité et permettent de modéliser des 
relations non linéaires, même avec des jeux de données de taille modérée. 
Elles reposent aussi sur des bases théoriques solides et peuvent être adaptées 
à différents types de données.


Elles présentent toutefois des limites : leur coût de calcul augmente fortement 
avec le nombre d'échantillons, leurs performances dépendent du réglage des 
hyperparamètres et leurs décisions sont parfois difficiles à interpréter. 
La prédiction peut aussi être lente, et l'obtention de probabilités nécessite 
une étape de calibrage supplémentaire.


Les SVM peuvent souffrir de surapprentissage lorsque leur modèle est trop complexe, ou de 
sous-apprentissage lorsqu'il est trop simple. Le réglage des hyperparamètres et la validation 
croisée permettent de trouver un bon équilibre et d'améliorer leur capacité à généraliser à de 
nouvelles données.

\cite{ref_01}

\newpage
\bibliography{report_v1.0}


\end{document}